% Author: Alfredo Sánchez Alberca (asalber@ceu.es)
% !TEX program = xelatex
\documentclass[a4paper, 10pt]{article}
  \usepackage[top=2.5cm, bottom=3cm, left=2cm, right=2cm, headsep=1.5cm]{geometry}
  %\usepackage{fontspec}
\usepackage{amsmath}
\usepackage{mathspec}
\usepackage[sfdefault,semibold,Light]{FiraSans}
\usepackage{newtxsf}
\usepackage{multicol}
\usepackage{tikz}
\usepackage{graphicx}
\usepackage[most]{tcolorbox}
\usepackage{fancyhdr}
\pagestyle{fancy}
\usepackage[pdfauthor={Alfredo Sánchez Alberca}, pdftitle={Statistical formulas}, colorlinks=true]{hyperref}

% COLORS
\usepackage{xcolor}
\definecolor{color1}{RGB}{5,161,230} % Blue
\definecolor{color2}{RGB}{238,50,36} % Red
\definecolor{color3}{RGB}{0,205,0} % Green
\definecolor{color4}{RGB}{243,102,25} % Orange

% LISTS
\usepackage[shortlabels]{enumitem} % Customize lists
%\setlist{nolistsep} % Reduce spacing between bullet points and numbered lists
\setlist[description]{style=sameline,leftmargin=0cm}

% TABLES
\usepackage{booktabs}

% GRAPHICS
\usepackage{pgfplots}
\usepgfplotslibrary{fillbetween}
% Line styles for plotting 2d functions
\pgfplotscreateplotcyclelist{2dfun cycle}{
{color1, line width=0.8pt},
{color2, line width=0.8pt},
{color3, line width=0.8pt},
{color4, line width=0.8pt},
}
\pgfplotsset{
  compat=1.14,
% Basic 2D settings
  2dbaseplot/.style={
    axis line style={<->}, % arrows on the axis
    scale only axis,        % otherwise width won't be as intended: http://tex.stackexchange.com/questions/36297/pgfplots-how-can-i-scale-to-text-width
    x tick label style={font=\small},
    y tick label style={font=\small},
    every axis/.append style={font=\small}, 
  },
% Settings for generic 2D functions 
  gen2dfun/.style={
    2dbaseplot,
    axis x line=middle,
    axis y line=middle,
    xlabel={$x$},         
    ylabel={$y$},          
    every axis x label/.style={at={(ticklabel* cs:1)}, anchor=west,},
	every axis y label/.style={at={(ticklabel* cs:1)}, anchor=south,},
    cycle list name=2dfun cycle,
  },
}
\usepackage{tikz}

% Sections
\usepackage{titlesec}
\titleformat{\section}
{\color{color1}\normalfont\Huge\bfseries}
{\color{color1}\thesection}{1em}{}
\titleformat{\subsection}
{\color{color1}\normalfont\Large\bfseries}
{\color{color1}\thesubsection}{1em}{}

\lhead{\includegraphics{img/logo-uspceu.pdf}}
\rhead{\url{http://aprendeconalf.es}}
\renewcommand{\headrulewidth}{0pt}
\renewcommand{\floatpagefraction}{.8}
\renewcommand{\textfraction}{.1}

\setlength{\columnsep}{1cm}

\DeclareMathOperator{\sen}{sen}
\DeclareMathOperator{\tg}{tg}
\DeclareMathOperator{\arcsen}{arcsen}
\DeclareMathOperator{\arctg}{arctg}
\DeclareMathOperator{\arcsec}{arcsec}
\DeclareMathOperator{\arccot}{arccot}
\DeclareMathOperator{\arccsc}{arccsc}

\begin{document}
% Author: Alfredo Sánchez Alberca (asalber@ceu.es)
\section*{Integrales}
\small
\tcbset{enhanced, colback=color1!10, colframe=color1, fonttitle=\bfseries\large\sffamily, %lifted shadow={1mm}{-2mm}{3mm}{0.1mm}
}

\begin{multicols}{2}
\begin{tcolorbox}[hbox, title=Primitivas básicas]
\begin{minipage}{0.4\textwidth}
\begin{description}
  \item[Potencias]
  \[
    \int a\,dx=ax+C
  \]
  \[
    \int x^n\,dx=\dfrac{x^{n+1}}{n+1}+C
  \]
  \item[Exponenciales]
  \[
    \int e^x\,dx=e^x+C.
  \]
  \[
    \dint a^x\,dx=\dfrac{a^x}{\ln(a)}+C
  \]
  \item[Logarítmos]
  \[
    \int \dfrac{1}{x}\, dx=\ln|x|+C
  \]
  \item[Trigonométricas]
  \[
    \int \sen(x)\, dx=-\cos(x)+C
  \]
  \[
    \int \cos(x)\, dx=\sen(x)+C
  \]
  \[
    \int \frac{1}{\cos(x)^2}\, dx = \tg(x)+C
  \]
  \[
    \int \tg(x)\, dx=\ln|\sec(x)|+C
  \]
  \[
    \int \sec(x)\, dx = \ln|\sec(x) + \tg(x)|+C
  \]
  \[
    \int \csc(x)\, dx= \ln|\csc(x)-\cot(x)|+C
  \] 
  \[
    \int \cot(x) \, dx= \ln|\sen(x)|+C
  \] 
  \[
    \int \sec(x)^2\, dx= \tg(x)+ C
  \] 
  \[
    \int \csc(x)^2\, dx= -\cot(x)+ C
  \] 
  \[
    \int \sec(x)\tg(x)\, dx= \sec(x)+ C
  \] 
  \[
    \int \csc(x)\cot(x)\, dx = -\csc(x) +C
  \]
  \item[Racionales e irracionales]
  \[
    \int\frac{dx}{\sqrt{a^2-x^2}}=\arcsen\left(\frac{x}{a}\right)+C
  \] 
  \[
    \int\frac{dx}{a^2+x^2}=\frac{1}{a}\arctg\left(\dfrac{x}{a}\right)+C
  \] 
  \[
    \int\dfrac{dx}{x\sqrt{x^2-a^2}}=\frac{1}{a}\arcsec\left(\dfrac{x}{a}\right)+C
  \] 
  \[
    \int\frac{dx}{a^2-x^2}=\frac{1}{2a}\ln\left|\frac{x+a}{x-a}\right|+C
  \]
\end{description}
\end{minipage}
\end{tcolorbox}

\begin{tcolorbox}[hbox, title=Reglas de integración]
  \begin{minipage}{0.4\textwidth}
    \begin{description}
      \item[Suma]
      \[
        \int u+v\,dx = \int u\,dx + \int v\,dx
      \]
      \item[Diferencia]
      \[
        \int u-v\,dx = \int u\,dx - \int v\,dx
      \]
      \item[Producto por una constante]
      \[
        \int af(x)\,dx = a\int f(x)\,dx
      \]
      \item[Partes]
      \[
        \int u\,dv = uv - \int v\,du
      \]
      \item[Sustitución]
      \[
        \int f(u)u'\,dx = \int f(u)\,du
      \]
    \end{description}
  \end{minipage}
\end{tcolorbox}

\begin{tcolorbox}[hbox, title=Integrales definidas]
  \begin{minipage}{0.4\textwidth}
    \begin{description}
      \item[Regla de Barrow]
      \[
        \int_a^b f(x)\,dx = F(b)-F(a)
      \]
      donde $F(x) = \int f(x)\,dx$
      \begin{center}
      % Author: Alfredo Sánchez Alberca (asalber@ceu.es)
\begin{tikzpicture}[trim axis left, trim axis right]
  \begin{axis}[
    gen2dfun, 
    xmin=0, xmax=4,
    ymin=0, ymax=4,
    xtick={1,3},
    xticklabels={$a$, $b$},
    ytick={},
    yticklabels={},
    clip=false, 
    height=3.5cm,
  	]
    \addplot+[domain=0.7:3.3, samples=100, smooth, thick, name path=F] {x^3-6*x^2+11*x-3} node[anchor=south west] {$f(x)$};
    \path[name path=axis] (axis cs:1,0) -- (axis cs:3,0);
    \addplot [thick, color=color1, fill=color1, fill opacity=0.3] fill between[of=F and axis, soft clip={domain=1:3},];
    \node at (2,1.5) {$\displaystyle\int_a^b f(x)\,dx$};
    \coordinate (O) at (0,0);
    \coordinate (A) at (1,3);
    \coordinate (B) at (3,3);
    \draw[gray, dashed] (A) -- (A|-O);
    \draw[gray, dashed] (B) -- (B|-O);
  \end{axis};
\end{tikzpicture}

      \end{center}
      \item[Propiedades]
      \[
        \int_a^a f(x)\,dx = 0
      \]
      \[
        \int_a^b f(x)\,dx = -\int_b^a f(x)\,dx
      \]
      \[
        \int_a^b f(x)\, dx = \int_a^c f(x)\,dx + \int_c^b f(x)\,dx
      \]
      donde $a<c<b$.
    \end{description}
  \end{minipage}
\end{tcolorbox}
\end{multicols}
\end{document}
